\documentclass[11pt,a4paper]{article}
\usepackage[T1]{fontenc}
\usepackage[utf8]{inputenc}
\usepackage{amsmath,amsthm,mathtools}
\usepackage{newtxtext,newtxmath}
\usepackage[margin=27mm,headheight=15pt]{geometry}
\usepackage{microtype,booktabs,array,longtable}
\usepackage{xcolor}
\usepackage{enumitem}
\usepackage{fancyhdr}
\usepackage{listings}
\usepackage{xurl}
\usepackage[colorlinks=true,linkcolor=blue!45!black,citecolor=blue!45!black,urlcolor=blue!45!black]{hyperref}
\usepackage{aliascnt}
\usepackage[nameinlink,noabbrev]{cleveref}
\hypersetup{pdftitle={New Deductions from Thue--Morse Cancellation},pdfauthor={Research article prepared for Vladimir Reshetnikov},pdfsubject={Spline corrections, dyadic reconstruction, local limit expansions, and spectral moment indeterminacy}}
\numberwithin{equation}{section}
\newtheorem{theorem}{Theorem}[section]
\newaliascnt{proposition}{theorem}
\newtheorem{proposition}[proposition]{Proposition}
\aliascntresetthe{proposition}
\newaliascnt{lemma}{theorem}
\newtheorem{lemma}[lemma]{Lemma}
\aliascntresetthe{lemma}
\newaliascnt{corollary}{theorem}
\newtheorem{corollary}[corollary]{Corollary}
\aliascntresetthe{corollary}
\theoremstyle{definition}
\newaliascnt{definition}{theorem}
\newtheorem{definition}[definition]{Definition}
\aliascntresetthe{definition}
\newaliascnt{example}{theorem}
\newtheorem{example}[example]{Example}
\aliascntresetthe{example}
\newaliascnt{remark}{theorem}
\newtheorem{remark}[remark]{Remark}
\aliascntresetthe{remark}
\newaliascnt{problem}{theorem}
\newtheorem{problem}[problem]{Research problem}
\aliascntresetthe{problem}
\crefname{proposition}{proposition}{propositions}
\Crefname{proposition}{Proposition}{Propositions}
\crefname{lemma}{lemma}{lemmas}
\Crefname{lemma}{Lemma}{Lemmas}
\crefname{corollary}{corollary}{corollaries}
\Crefname{corollary}{Corollary}{Corollaries}
\crefname{definition}{definition}{definitions}
\Crefname{definition}{Definition}{Definitions}
\crefname{example}{example}{examples}
\Crefname{example}{Example}{Examples}
\crefname{remark}{remark}{remarks}
\Crefname{remark}{Remark}{Remarks}
\crefname{problem}{research problem}{research problems}
\Crefname{problem}{Research problem}{Research problems}
\newcommand{\R}{\mathbb R}
\newcommand{\N}{\mathbb N}
\newcommand{\Z}{\mathbb Z}
\newcommand{\E}{\mathbb E}
\newcommand{\Prob}{\mathbb P}
\newcommand{\eps}{\varepsilon}
\newcommand{\dd}{\,\mathrm d}
\newcommand{\ii}{\mathrm i}
\newcommand{\D}{\mathrm D}
\newcommand{\sinc}{\operatorname{sinc}}
\newcommand{\supp}{\operatorname{supp}}
\newcommand{\Res}{\operatorname{Res}}
\newcommand{\norm}[1]{\left\lVert #1\right\rVert}
\newcommand{\abs}[1]{\left\lvert #1\right\rvert}
\newcommand{\Law}{\mathcal L}
\newcommand{\Part}[1]{\substack{k_1,\ldots,k_{#1}\geq0\\\sum_{r=1}^{#1}rk_r=#1}}
\newcommand{\QP}[2]{(#1;#2)}
\DeclareMathOperator{\dist}{dist}
\pagestyle{fancy}
\fancyhf{}
\fancyhead[L]{\small Thue--Morse cancellation: new deductions}
\fancyhead[R]{\small Research article}
\fancyfoot[C]{\thepage}
\setlength{\parindent}{1.2em}
\setlength{\parskip}{3pt}
\setlist{itemsep=3pt,topsep=5pt}
\lstset{basicstyle=\ttfamily\footnotesize,breaklines=true,columns=fullflexible,frame=single,framerule=0.3pt,showstringspaces=false}
\emergencystretch=2em

\begin{document}
\begin{titlepage}
\centering
\vspace*{12mm}
{\large A research extension of the Thue--Morse Atlas\par}
\vspace{15mm}
{\LARGE\bfseries New Deductions from\\[4pt] Thue--Morse Cancellation\par}
\vspace{8mm}
{\Large Global spline corrections, exact dyadic reconstruction,\\[3pt]
local-limit expansions, and spectral moment twins\par}
\vspace{11mm}
{\large 4 September 2026\par}
\vspace{12mm}
\begin{minipage}{0.91\textwidth}
\begin{abstract}
The factorization of a Thue--Morse block into binary differences connects
integer cancellation, compactly supported smooth functions, and lacunary
Fourier products. We exploit the geometry of the derivative copies, rather
than only the formal product expansion, to obtain four groups of results.
First, finite uniform-convolution splines admit a global, all-orders
$C^k$ correction expansion with explicit error bounds and exact leading
supremum constants; no neighborhoods of their knots must be removed.
Second, at a fixed dyadic argument the finite spline values become an
\emph{exact polynomial} in $4^{-m}$. A finite Richardson formula therefore
recovers the limiting value exactly, with absolute error amplification less
than two at every extrapolation order.
Third, the positive coefficients obtained by repeatedly summing a finite
Thue--Morse block satisfy an exact differential conversion formula, a
uniform local-limit expansion with explicit coefficient polynomials, and
an exact terminating version on dyadic lattice subsequences. Finally, the
normalized Fourier-energy density has moments $2^{r(r+1)}$ and consequently
shares every polynomial moment with both a symmetric lognormal density and
an explicit family of discrete theta measures. A separated geometric
extension explains this last phenomenon beyond the binary case.

All central claims are proved. Exact rational computations are supplied
separately for reproducibility. The claimed contribution is the collection
of deductions and quantitative refinements established here; classical
identities are identified, and worldwide publication priority is not
asserted.
\end{abstract}
\end{minipage}
\vfill
\begin{minipage}{0.91\textwidth}
\small
\textbf{Status.} An independently checkable research manuscript, not a
peer-reviewed publication or a Lean formalization. It develops and tests
specific extensions of the supplied repository document. The distinction
between exact identities, asymptotic statements, and unresolved extensions
is maintained throughout.
\end{minipage}
\end{titlepage}
\begingroup
\small
\setlength{\parskip}{0pt}
\tableofcontents
\endgroup
\newpage

\section{Research setting and the main contributions}\label{sec:intro}

The signed Thue--Morse sequence is
\begin{equation}
 \eps_n=(-1)^{s_2(n)},\qquad n\geq0,
 \label{eq:epsilon}
\end{equation}
where $s_2(n)$ is the number of ones in the binary expansion of $n$.
Its apparent irregularity coexists with exact cancellation on every
complete binary block. The classical survey of Allouche and Shallit
\cite{allouche-shallit} describes the unusually broad reach of this
sequence. Here we investigate one particularly productive direction:
what quantitative information survives when those finite cancellations
are integrated repeatedly and then passed to a smooth limit?

The supplied \emph{Thue--Morse Atlas and Frontiers} \cite{atlas} already
contains the finite-block factorization, the Fabius--Rvachev bridge,
spectral correction coefficients, and a weak spline expansion. Its
Part~II leaves a strong spline-correction statement away from knots as a
conjecture. Our first main theorem proves a global fixed-order statement,
including the knots, and identifies the exact leading norm constant.
The other principal developments concern finite stopping on dyadic
arguments, the local behavior of repeated-prefix coefficients, and a
moment-theoretic interpretation of Fourier energy.

The underlying compactly supported function is not new. Its probability
interpretation, derivative identities, and rational values at dyadic
points belong to the established theory, particularly the work of Arias
de Reyna \cite{arias1982,arias2018}. We use these mechanisms explicitly,
reprove the required facts, and do not relabel rationality or the binary
product identity as discoveries. Recent work on Thue--Morse transforms,
for example \cite{cloitre}, provides additional context for finite-block
identities; the analytic questions treated here are more specific.

\subsection{A compact statement of the results}

Let $V_1,V_2,\ldots$ be independent, uniformly distributed random variables
on $[-1,1]$. Let $u_m$ and $u$ be the densities of
\begin{equation}
 Y_m=\sum_{j=1}^m2^{-j}V_j,
 \qquad
 Y=\sum_{j=1}^{\infty}2^{-j}V_j,
 \qquad h=2^{-m}.
 \label{eq:random-series}
\end{equation}
Thus $u$ is the Rvachev up-function in the normalization
$\supp u=[-1,1]$ and $u(0)=1$. Write
\begin{equation}
 a_j=\frac{\E Y^{2j}}{(2j)!},\qquad
 K_r=2^{r(r+1)/2}.
 \label{eq:a-K}
\end{equation}
The four main conclusions are as follows.

\paragraph{Global correction with a sharp constant.}
For fixed $k,N\geq0$, and $m\geq k+2N+4$,
\begin{equation}
 \left\|\D^k\left(u-\sum_{j=0}^N a_j4^{-jm}u_m^{(2j)}\right)\right\|_\infty
 \leq a_{N+1}K_{k+2N+2}4^{-(N+1)m}.
 \label{eq:intro-strong}
\end{equation}
After multiplication by $4^{(N+1)m}$, the norm on the left converges to
the constant on the right. In particular,
$4^m\norm{u-u_m}_\infty\to4/9$.
The proof is in \cref{sec:global}.

\paragraph{Exact dyadic reconstruction.}
If $x=a/2^d\in(-1,1)$, then, once $m\geq d+2$,
\begin{equation}
 u_m(x)=\sum_{j=0}^{\lfloor d/2\rfloor}b_j4^{-jm}u^{(2j)}(x),
 \qquad \sum_{j\geq0}b_jz^{2j}=\frac1{\E e^{zY}}.
 \label{eq:intro-dyadic}
\end{equation}
This is an equality, not an asymptotic series. Consequently
$\lfloor d/2\rfloor+1$ consecutive spline levels determine $u(x)$ exactly.
The reconstruction weights have total absolute value less than two.
See \cref{sec:dyadic}.

\paragraph{A local-limit expansion with finite stopping.}
Let $c_m(k)$ be the coefficients of
\begin{equation}
 A_m(z)=\frac{\prod_{j=0}^{m-1}(1-z^{2^j})}{(1-z)^m}
       =\prod_{j=0}^{m-1}(1+z+\cdots+z^{2^j-1}).
 \label{eq:intro-Am}
\end{equation}
After the centering and scaling in \cref{sec:lattice}, their lattice
density $\rho_m$ satisfies uniformly on its entire lattice
\begin{equation}
 \rho_m(x)=u(x)-\left(\frac m6+\frac1{18}\right)4^{-m}u''(x)
            +O(m^2 16^{-m}).
 \label{eq:intro-local}
\end{equation}
There are explicit polynomial coefficients at every order. At fixed
dyadic points, on the odd-level lattice subsequence, the expansion
terminates exactly. The sharp uncorrected uniform error constant is
$4/3$ when the error is divided by $m4^{-m}$.

\paragraph{Moment twins of Fourier energy.}
With angular-frequency Fourier transform
$Q(t)=\int e^{\ii tx}u(x)\dd x$, the probability measure
\begin{equation}
 \dd\nu(t)=\frac{|Q(t)|^2\dd t}{\int_{\R}|Q(s)|^2\dd s}
 \label{eq:intro-nu}
\end{equation}
has odd moments zero and even moments $2^{r(r+1)}$. It has the same
polynomial moments as $S e^L$, where $S$ is an independent uniform sign
and
\begin{equation}
 L\sim N\left(\frac{\log2}{2},\frac{\log2}{2}\right).
 \label{eq:intro-lognormal}
\end{equation}
An explicit one-parameter family of purely atomic measures has these
moments as well. The identification of this spectral measure, rather
than the classical indeterminacy of lognormal moments itself, is the
point of \cref{sec:spectral}.

\subsection{Scope of the novelty claim}

There are two different questions: whether a theorem is proved, and
whether it has appeared anywhere before. The proofs below settle the
first for the stated results. A targeted literature check supports
presenting the global sharp correction theorem, finite dyadic stopping,
and the explicit lattice refinements as research contributions relative
to the inspected atlas. That check is not an exhaustive priority search.
The Fourier-energy identification is presented as an additional
potentially novel connection, with the lognormal background explicitly
credited to the classical moment literature \cite{heyde}.

Nothing below purports to solve questions about primes, squares, cubes,
normality, or all possible Thue--Morse exponential sums. We instead
establish a coherent set of complete results in the finite-block,
spline, and Fourier directions.

\section{From binary cancellation to compactly supported splines}\label{sec:basic}

\subsection{The finite product and polynomial cancellation}

For $m\geq0$ define
\begin{equation}
 P_m(z)=\sum_{n=0}^{2^m-1}\eps_nz^n.
 \label{eq:Pm}
\end{equation}
Choosing whether each binary digit is present gives the classical identity
\begin{equation}
 P_m(z)=\prod_{j=0}^{m-1}(1-z^{2^j}).
 \label{eq:block-product}
\end{equation}
Indeed, a subset $J\subseteq\{0,\ldots,m-1\}$ contributes
$(-1)^{|J|}z^{\sum_{j\in J}2^j}$, and every integer in the block has a
unique such representation.

Substitution $z=e^t$ gives
\begin{equation}
 \sum_{n=0}^{2^m-1}\eps_ne^{nt}
 =(-1)^m2^{m(m-1)/2}t^m+O(t^{m+1}),
\end{equation}
so
\begin{equation}
 \sum_{n=0}^{2^m-1}\eps_n n^r=
 \begin{cases}
 0,&0\leq r<m,\\
 (-1)^m m!2^{m(m-1)/2},&r=m.
 \end{cases}
 \label{eq:prouhet}
\end{equation}
This elementary Prouhet cancellation is the algebraic source of all
compact-support cancellations used below.

\subsection{The continuous convolution model}

The $j$th term $2^{-j}V_j$ has density
\begin{equation}
 f_j(x)=2^{j-1}\mathbf1_{[-2^{-j},2^{-j}]}(x).
\end{equation}
Therefore $u_m=f_1*\cdots*f_m$. Its support is
$[-1+h,1-h]$, and for $m\geq2$ it is a $C^{m-2}$ piecewise polynomial
of degree at most $m-1$.

The infinite random series in \eqref{eq:random-series} converges
absolutely for every outcome. The distributional tail is an exact
scaled copy of the entire series:
\begin{equation}
 Y\ \overset{\mathrm d}=\ Y_m+hY',
 \qquad Y'\text{ independent of }Y_m,\quad \Law(Y')=\Law(Y).
 \label{eq:tail-law}
\end{equation}
Its characteristic function is
\begin{equation}
 Q(t)=\prod_{j=1}^{\infty}\sinc(2^{-j}t),\qquad
 \sinc t=\begin{cases}\sin t/t,&t\ne0,\\1,&t=0.\end{cases}
 \label{eq:Q}
\end{equation}
Our Fourier convention is
\begin{equation}
 Q(t)=\int_{\R}e^{\ii tx}u(x)\dd x,
 \qquad
 u(x)=\frac1{2\pi}\int_{\R}e^{-\ii tx}Q(t)\dd t.
 \label{eq:Fourier-convention}
\end{equation}
No factor of $2\pi$ is hidden inside the definition of $\sinc$.

\begin{lemma}[Smoothness and a useful Fourier bound]\label{lem:decay}
For every integer $L\geq1$,
\begin{equation}
 |Q(t)|\leq\min\{1,K_L|t|^{-L}\},
 \qquad K_L=2^{L(L+1)/2}.
 \label{eq:Qbound}
\end{equation}
The law of $Y$ has an even density $u\in C_c^\infty(\R)$, supported
on $[-1,1]$, with $0\leq u\leq1$ and $u(0)=1$.
\end{lemma}
\begin{proof}
For the first $L$ factors, use
$|\sinc(2^{-j}t)|\leq2^j/|t|$; all remaining factors have modulus
at most one. This proves \eqref{eq:Qbound}. Consequently
$|t|^rQ(t)$ is integrable for every $r$, by choosing $L>r+1$.
Fourier inversion yields a smooth density. Compact support and evenness
follow from the random series. Finally write $Y=V_1/2+R$, where
$R$ is independent and supported on $[-1/2,1/2]$. Convolution with the
unit-height density of $V_1/2$ proves $u\leq1$ and $u(0)=1$.
The same reasoning shows $\norm{u_m}_\infty=1$.
\end{proof}

\subsection{A finite signed formula}

For $y\in\R$ write $y_+=\max\{y,0\}$.
\begin{proposition}[Thue--Morse positive-part formula]\label{prop:positive-part}
For $m\geq2$,
\begin{equation}
 u_m(x)=\frac{2^{m(m-1)/2}}{(m-1)!}
 \sum_{n=0}^{2^m-1}\eps_n
       (x+1-h-2hn)_+^{m-1}.
 \label{eq:spline-positive}
\end{equation}
In particular the knots are among $-1+h+2hn$.
For $0\leq r\leq m-2$, differentiation gives
\begin{equation}
 u_m^{(r)}(x)=\frac{2^{m(m-1)/2}}{(m-1-r)!}
 \sum_{n=0}^{2^m-1}\eps_n
       (x+1-h-2hn)_+^{m-1-r}.
 \label{eq:spline-deriv-positive}
\end{equation}
The analogous formula for $r=m-1$ holds on open cells.
\end{proposition}
\begin{proof}
The distributional derivative of $f_j$ is
$2^{j-1}(\delta_{-2^{-j}}-\delta_{2^{-j}})$. Hence
\begin{equation}
 u_m^{(m)}=2^{m(m-1)/2}
 \sum_{n=0}^{2^m-1}\eps_n\delta_{-1+h+2hn}.
\end{equation}
To see the locations and signs, select an upper endpoint in any subset
of the $m$ factors. The lengths $2^{1-j}$, when put in reverse order,
are $2h,4h,\ldots,2^mh$. Their subset sums give $2hn$ and their signs
are $\eps_n$. Integrating $m$ times from $-\infty$ proves
\eqref{eq:spline-positive}. The polynomial terms to the right of the
support cancel by \eqref{eq:prouhet}. The derivative formula follows.
\end{proof}

The signed formula is excellent for exact rational computation, but
its large cancelling summands make naive floating-point evaluation
unreliable. A stability statement for extrapolation is not a stability
statement for this inner signed sum.

\section{The derivative-copy geometry}\label{sec:derivative}

The key analytic fact is that, after differentiation, the rescaled
copies have disjoint interiors. Thus differentiating a fixed number of
times does \emph{not} introduce a factor growing with $m$.

\begin{proposition}[Exact derivative copies and norm identities]\label{prop:copies}
For integers $r\geq0$ and $m\geq r+2$,
\begin{align}
 u_m^{(r)}(x)
 &=K_r\sum_{n=0}^{2^r-1}\eps_n
       u_{m-r}(2^r x+2^r-1-2n),\label{eq:finite-copies}\\
 u^{(r)}(x)
 &=K_r\sum_{n=0}^{2^r-1}\eps_n
       u(2^r x+2^r-1-2n).\label{eq:infinite-copies}
\end{align}
For every $1\leq p\leq\infty$,
\begin{align}
 \norm{u_m^{(r)}}_p&=K_r\norm{u_{m-r}}_p,
 &\norm{u^{(r)}}_p&=K_r\norm{u}_p,\label{eq:Lp-copies}\\
 \norm{u_m^{(r)}}_\infty&=K_r,
 &\norm{u^{(r)}}_\infty&=K_r.\label{eq:sharp-deriv-norm}
\end{align}
\end{proposition}
\begin{proof}
Differentiating the first box and rescaling the tail gives
\begin{equation}
 u_m'(x)=2\{u_{m-1}(2x+1)-u_{m-1}(2x-1)\}.
 \label{eq:refinement-derivative}
\end{equation}
The same argument applies to $u$. On iterating, the accumulated
factor is $2^{1+2+\cdots+r}=K_r$. At every step, endpoint choices
contribute binary digit signs, proving the two copy identities.

The $n$th infinite copy is supported on
\begin{equation}
 \left[-1+\frac{2n}{2^r},\ -1+\frac{2n+2}{2^r}\right].
\end{equation}
These intervals meet only at their endpoints. Finite copies have still
smaller supports. Thus there is no cancellation in the integral of the
$p$th power of the absolute value. A change of variable contributes
$2^{-r}$ for each of the $2^r$ copies, giving
\eqref{eq:Lp-copies}. The supremum statement uses
$\norm{u_j}_\infty=\norm{u}_\infty=1$.
\end{proof}

\begin{remark}[What is classical, and what is used here]
Derivative self-similarity is an established part of the Fabius theory
\cite{arias1982,arias2018}. We emphasize the exact norm consequence and
its uniformity in the finite level. That simple observation is the
mechanism behind the global remainder theorem and the spectral-moment
calculation below.
\end{remark}

\begin{corollary}[Vanishing high derivatives at a dyadic point]\label{cor:dyadic-jet}
If $x=a/2^d\in(-1,1)$ and $r>d$, then $u^{(r)}(x)=0$.
\end{corollary}
\begin{proof}
In \eqref{eq:infinite-copies}, $2^rx$ is an even integer and
$2^r-1-2n$ is odd. Every argument is therefore an odd integer.
The density $u$ is zero at and outside $\pm1$.
\end{proof}

This does not say that $u$ agrees with its Taylor polynomial on a
neighborhood of $x$. Vanishing higher derivatives at one point and
local polynomiality are very different assertions for a $C^\infty$
function.

\section{Explicit coefficient calculus}\label{sec:coefficients}

The moment generating function
\begin{equation}
 M(z)=\E e^{zY}=\prod_{j=1}^{\infty}
          \frac{\sinh(2^{-j}z)}{2^{-j}z}
      =\sum_{j=0}^{\infty}a_jz^{2j}
 \label{eq:MGF}
\end{equation}
is entire. Its reciprocal is analytic near the origin. We shall need
both its coefficients and those of two closely related functions.

Define
\begin{equation}
 \lambda_r=\frac{\zeta(2r)}{r\pi^{2r}},\qquad r\geq1.
 \label{eq:lambda}
\end{equation}
The sine product \cite[Sec.~4.22]{dlmf}, or its logarithmic derivative
followed by integration, gives
\begin{align}
 -\log\sinc z&=\sum_{r\geq1}\lambda_rz^{2r},\label{eq:log-sinc}\\
 \log M(z)&=\sum_{r\geq1}
     \frac{(-1)^{r+1}\lambda_r}{4^r-1}z^{2r}.
 \label{eq:logM}
\end{align}
The first coefficients are
$\lambda_1=1/6$, $\lambda_2=1/180$, and $\lambda_3=1/2835$.
Euler's even-zeta formula makes every $\lambda_r$ rational.

For any sequence $v_1,v_2,\ldots$, define the finite partition polynomial
\begin{equation}
 \mathcal E_j(v)=
 \sum_{\Part{j}}\prod_{r=1}^j\frac{v_r^{k_r}}{k_r!},
 \qquad \mathcal E_0(v)=1.
 \label{eq:partition-polynomial}
\end{equation}
It is characterized by
$\exp(\sum_{r\geq1}v_rz^r)=\sum_{j\geq0}\mathcal E_j(v)z^j$.
Thus all coefficients in this article have explicit finite-sum
expressions, without requiring auxiliary recurrences.

\begin{proposition}[The four coefficient families]\label{prop:coeff-families}
Let
\begin{align}
 \frac1{M(z)}&=\sum_{j\geq0}b_jz^{2j},\label{eq:b-def}\\
 \left(\frac z{\sinh z}\right)^m&=
      \sum_{j\geq0}\gamma_j(m)z^{2j},\label{eq:gamma-def}\\
 \exp\left(\sum_{r\geq1}\lambda_r
           \left(m+\frac1{4^r-1}\right)z^r\right)
      &=\sum_{j\geq0}H_j(m)z^j.\label{eq:H-def}
\end{align}
Then
\begin{align}
 a_j&=\mathcal E_j\left(
        \frac{(-1)^{r+1}\lambda_r}{4^r-1}\right),\label{eq:a-partition}\\
 b_j&=(-1)^j\mathcal E_j\left(\frac{\lambda_r}{4^r-1}\right),\label{eq:b-partition}\\
 \gamma_j(m)&=(-1)^j\mathcal E_j(m\lambda_r),\label{eq:gamma-partition}\\
 H_j(m)&=\mathcal E_j\left(\lambda_r
                     \left(m+\frac1{4^r-1}\right)\right).
 \label{eq:H-partition}
\end{align}
For each $j$, $H_j(m)$ is a rational polynomial of degree $j$ with
positive coefficients and leading term $m^j/(6^j j!)$. Moreover,
\begin{equation}
 \sum_{\ell=0}^j\gamma_\ell(m)b_{j-\ell}=(-1)^jH_j(m).
 \label{eq:coefficient-convolution}
\end{equation}
\end{proposition}
\begin{proof}
Expand the exponentials in \eqref{eq:logM}--\eqref{eq:H-def}.
For the alternating families, the sign of each partition monomial
is $(-1)^{\sum rk_r}=(-1)^j$. The only degree-$j$ term in $m$ comes
from $k_1=j$. Multiplication of the two reciprocal generating
functions gives \eqref{eq:coefficient-convolution}.
\end{proof}

For reference, the first $a_j$ and $b_j$ are
\begin{center}
\renewcommand{\arraystretch}{1.45}
\begin{tabular}{c@{\qquad}r@{\qquad}r}
\toprule
$j$ & $a_j$ & $b_j$\\
\midrule
0 & $1$ & $1$\\
1 & $1/18$ & $-1/18$\\
2 & $19/16200$ & $31/16200$\\
3 & $583/42865200$ & $-2347/42865200$\\
4 & $132809/1311675120000$ & $1904369/1311675120000$\\
\bottomrule
\end{tabular}
\end{center}
The first lattice correction polynomials are
\begin{align}
 H_0(m)&=1,\nonumber\\
 H_1(m)&=\frac m6+\frac1{18},\label{eq:H-first}\\
 H_2(m)&=\frac{(m+1/3)^2}{72}+\frac{m+1/15}{180},\nonumber\\
 H_3(m)&=\frac{(m+1/3)^3}{1296}
      +\frac{(m+1/3)(m+1/15)}{1080}
      +\frac{m+1/63}{2835}.\nonumber
\end{align}

\subsection{An optional rational recurrence for computation}

Writing $\mu_{2j}=\E Y^{2j}$, one may alternatively use
\begin{equation}
 (2j+1)(4^j-1)\mu_{2j}
 =\sum_{\ell=0}^{j-1}\binom{2j+1}{2\ell}\mu_{2\ell},
 \qquad \mu_0=1.
 \label{eq:moment-recurrence}
\end{equation}
For completeness, $M(2z)=M(z)\sinh z/z$. Multiplying by $z$,
comparing the coefficient of $z^{2j+1}$, and moving the last term
to the left proves \eqref{eq:moment-recurrence}. This recurrence is
convenient for exact arithmetic, but \eqref{eq:a-partition} is a
nonrecursive alternative.

\section{Global spline corrections, sharpness, and increasing order}\label{sec:global}

Define the corrected spline
\begin{equation}
 C_{m,N}(x)=\sum_{j=0}^N a_jh^{2j}u_m^{(2j)}(x).
 \label{eq:corrected-spline}
\end{equation}

\begin{theorem}[Global all-orders correction]\label{thm:global}
Let $k,N\geq0$ be integers, put $r=k+2N+2$, and suppose $m\geq r+2$.
Then
\begin{equation}
 \boxed{\ \norm{\D^k(u-C_{m,N})}_\infty
 \leq a_{N+1}K_rh^{2N+2}.\ }
 \label{eq:global-bound}
\end{equation}
The same right-hand side bounds
$\norm{u-C_{m,N}}_{C^k}:=\max_{0\leq\ell\leq k}
\norm{\D^\ell(u-C_{m,N})}_\infty$.
The estimate is global on $\R$, including all spline knots and the
endpoints of the supports.
\end{theorem}
\begin{proof}
The exact tail identity \eqref{eq:tail-law} gives
\begin{equation}
 u^{(k)}(x)=\E\,u_m^{(k)}(x-hY).
 \label{eq:tail-differentiation}
\end{equation}
Set $g=u_m^{(k)}$. By the hypothesis on $m$, $g$ has $2N+2$
continuous derivatives. Taylor's formula with the integral remainder
implies, for real $y$,
\begin{equation}
 \left|g(x-hy)-\sum_{\ell=0}^{2N+1}
             \frac{(-hy)^\ell}{\ell!}g^{(\ell)}(x)\right|
 \leq \frac{h^{2N+2}|y|^{2N+2}}{(2N+2)!}
             \norm{g^{(2N+2)}}_\infty.
\end{equation}
Average over $Y$. All odd powers disappear because $Y$ is symmetric.
The remaining finite sum is precisely $\D^k C_{m,N}$, and the
remainder is bounded by
$a_{N+1}h^{2N+2}\norm{u_m^{(r)}}_\infty$.
Now apply \eqref{eq:sharp-deriv-norm}. For lower derivatives use
$K_{\ell+2N+2}\leq K_r$.
\end{proof}

\begin{corollary}[Global H\"older control]\label{cor:holder}
Let $0<\alpha\leq1$, $r=k+2N+2$, and $m\geq r+3$. With
$[f]_\alpha=\sup_{x\ne y}|f(x)-f(y)|/|x-y|^\alpha$,
\begin{equation}
 [\D^k(u-C_{m,N})]_\alpha
 \leq 2^{1+\alpha r}a_{N+1}K_rh^{2N+2}.
 \label{eq:holder-bound}
\end{equation}
Thus a global fixed-order H\"older version is available as well.
\end{corollary}
\begin{proof}
For a bounded differentiable function, the two estimates
$|f(x)-f(y)|\leq2\norm f_\infty$ and
$|f(x)-f(y)|\leq\norm{f'}_\infty|x-y|$ imply
$[f]_\alpha\leq(2\norm f_\infty)^{1-\alpha}
\norm{f'}_\infty^\alpha$. Apply \cref{thm:global} at derivative
orders $k$ and $k+1$, and use $K_{r+1}=2^{r+1}K_r$.
\end{proof}

\paragraph{Why the knots are harmless here.}
The differentiation order is fixed before the level $m$ tends to
infinity. The spline degree and its classical smoothness grow with
$m$, while the norm of each admissible fixed derivative is exactly
constant in $m$. Thus there is no knot singularity in the derivatives
used in the Taylor remainder. A different argument would be required
for arbitrary derivative orders beyond the stated regularity threshold.

\begin{theorem}[Uniform leading remainder and sharp norm constant]\label{thm:sharp}
For fixed $k,N\geq0$,
\begin{equation}
 h^{-2N-2}\D^k(u-C_{m,N})
 \longrightarrow a_{N+1}u^{(k+2N+2)}
 \quad\text{uniformly on }\R.
 \label{eq:leading-remainder}
\end{equation}
Consequently
\begin{equation}
 \lim_{m\to\infty}4^{(N+1)m}
       \norm{\D^k(u-C_{m,N})}_\infty
       =a_{N+1}K_{k+2N+2}.
 \label{eq:sharp-limit}
\end{equation}
\end{theorem}
\begin{proof}
Apply \cref{thm:global} with order $N+1$ to obtain
\begin{equation}
 \D^k(u-C_{m,N})
 =a_{N+1}h^{2N+2}u_m^{(k+2N+2)}+O_{k,N}(h^{2N+4})
\end{equation}
in the supremum norm. The same theorem at correction order zero,
with derivative order $k+2N+2$, proves convergence of the displayed
spline derivative to $u^{(k+2N+2)}$. Uniform convergence implies
convergence of the supremum norms, and \cref{prop:copies} evaluates
the limiting norm.
\end{proof}

For the first two un-differentiated cases,
\begin{align}
 \lim_{m\to\infty}4^m\norm{u-u_m}_\infty&=\frac49,\label{eq:sharp-first}\\
 \lim_{m\to\infty}16^m
 \left\|u-u_m-\frac{4^{-m}}{18}u_m''\right\|_\infty
 &=\frac{2432}{2025}.
 \label{eq:sharp-second}
\end{align}
These constants are not estimates fitted to numerical data.
They are forced by the exact derivative norms.

\subsection{A certified increasing-order construction}

The next result allows the correction order to grow with the spline
level. It is an absolute error statement, not a relative endpoint
asymptotic.

\begin{corollary}[Supergeometric absolute accuracy]\label{cor:growing}
Fix $k\geq0$, let $m\geq k+4$, and choose $r$ to be the largest integer
with $r\leq m-2$ and $r\equiv k\pmod2$. Set
$N=(r-k-2)/2$. Then
\begin{equation}
 \norm{\D^k(u-C_{m,N})}_\infty
 \leq \frac{2^{-m^2/2+(k+1/2)m+3}}{(m-k-3)!}.
 \label{eq:growing-bound}
\end{equation}
\end{corollary}
\begin{proof}
Since $|Y|\leq1$, $a_{N+1}\leq1/(r-k)!$. The bound in
\cref{thm:global} becomes
\begin{equation}
 \frac{2^{r(r+1)/2-(r-k)m}}{(r-k)!}.
\end{equation}
Write $r=m-\delta$, where $\delta\in\{2,3\}$. The exponent equals
$-m^2/2+(k+1/2)m+\delta(\delta-1)/2$, and
$(r-k)!\geq(m-k-3)!$. This proves the claim.
\end{proof}

The large positive and negative terms in an actual evaluation still
need controlled arithmetic. The error certificate establishes the
quality of the exact corrected expression; it does not provide a
complete floating-point or bit-complexity analysis.

\section{Exact stabilization on dyadic arguments}\label{sec:dyadic}

Fix $x=a/2^d\in(-1,1)$. It is harmless to choose a nonminimal $d$;
a smaller denominator exponent simply gives a sharper degree bound.

\begin{lemma}[A low-degree cell around a dyadic point]\label{lem:dyadic-cell}
For $m\geq d+2$, the restriction of $u_m$ to
$[x-h,x+h]$ agrees with one polynomial of degree at most $d$.
\end{lemma}
\begin{proof}
The knots have coordinates $-1+h+2hn$, which are odd multiples of
$h$ when $m\geq1$. Since $2^mx$ is even, $x$ is the center of an
open cell whose endpoints are $x-h$ and $x+h$.

Apply the derivative-copy formula at order $d+1$ on this open cell.
It remains valid there in the borderline case as a piecewise derivative.
Put $\ell=m-d-1\geq1$. Each argument of $u_\ell$ has the form
\begin{equation}
 \text{an odd integer}+2^{d+1}(y-x),\qquad |y-x|<h.
\end{equation}
Its distance from zero is larger than $1-2^{-\ell}$, unless it is
farther away still. But $u_\ell$ is supported on
$[-1+2^{-\ell},1-2^{-\ell}]$. Thus $u_m^{(d+1)}=0$ on the cell.
The spline is a polynomial there, so its degree is at most $d$.
Continuity gives the endpoint equality as well.
\end{proof}

\begin{theorem}[Exact finite tail calculus at dyadic points]\label{thm:dyadic}
For $x=a/2^d\in(-1,1)$, $m\geq d+2$, and $0\leq r\leq d$,
\begin{align}
 u^{(r)}(x)&=\sum_{j=0}^{\lfloor(d-r)/2\rfloor}
           a_jh^{2j}u_m^{(r+2j)}(x),\label{eq:dyadic-forward}\\
 u_m^{(r)}(x)&=\sum_{j=0}^{\lfloor(d-r)/2\rfloor}
           b_jh^{2j}u^{(r+2j)}(x).\label{eq:dyadic-backward}
\end{align}
Both formulas are exact.
\end{theorem}
\begin{proof}
The tail formula \eqref{eq:tail-differentiation} averages the derivative
of $u_m$ at points $x-hY\in[x-h,x+h]$. On this cell the derivative
is a polynomial of degree at most $d-r$. Its Taylor formula has no
remainder. Symmetry gives \eqref{eq:dyadic-forward}.

View the vector of derivatives of orders $r,r+2,\ldots$ at $x$ as a
finite vector. The first formula is a triangular transformation with
diagonal entries one. The power-series identity $M(z)/M(z)=1$
shows that its inverse has coefficients $b_j$, which proves
\eqref{eq:dyadic-backward}. This is inversion of a finite triangular
system, not an assertion that an infinite differential operator
converges on an arbitrary smooth function.
\end{proof}

\begin{corollary}[Polynomial stabilization]\label{cor:poly-stabilize}
For fixed $x=a/2^d$, there is a polynomial $R_x$ of degree at most
$p=\lfloor d/2\rfloor$ such that
\begin{equation}
 u_m(x)=R_x(4^{-m}),\quad m\geq d+2,
 \qquad R_x(0)=u(x).
 \label{eq:exact-poly}
\end{equation}
Its coefficient of $z^j$ is $b_ju^{(2j)}(x)$.
\end{corollary}

\begin{corollary}[The degree is exact in lowest terms]\label{cor:exact-degree}
Suppose $d\geq1$ and $a$ is odd, so $a/2^d$ is in lowest terms.
The polynomial $R_x$ has degree exactly $p=\lfloor d/2\rfloor$.
More explicitly, its highest derivative factor is
\begin{equation}
 u^{(2p)}(a/2^d)=
 \begin{cases}
 K_d\eps_{(a+2^d-1)/2},&d\text{ even},\\[3pt]
 \frac12K_{d-1}\eps_{\lfloor(a+2^d)/4\rfloor},&d\text{ odd}.
 \end{cases}
 \label{eq:highest-dyadic-derivative}
\end{equation}
\end{corollary}
\begin{proof}
For even $d$, precisely one argument in the order-$d$ copy formula
is zero, and all the others are outside the support. For odd $d$,
at order $d-1$ precisely one argument is $1/2$ or $-1/2$.
The corresponding values are $u(0)=1$ and $u(\pm1/2)=1/2$.
The latter equality follows directly by convolving the first box
with its symmetric tail. The indices in
\eqref{eq:highest-dyadic-derivative} identify these copies. Both
values are nonzero. Finally $b_p\ne0$ by
\eqref{eq:b-partition}, so the highest coefficient of $R_x$
is nonzero.
\end{proof}

The crucial word is ``polynomial.'' An ordinary all-orders asymptotic
expansion, by itself, would not justify finite exact reconstruction.
Here the geometry of the dyadic cell supplies the missing exactness.

\subsection{Finite Richardson reconstruction}

For $0<q<1$ and $p\geq0$, let
\begin{equation}
 (q;q)_j=\prod_{\ell=1}^j(1-q^\ell),\qquad(q;q)_0=1,
\end{equation}
and define
\begin{equation}
 w_{p,j}(q)=
 \prod_{\substack{0\leq\ell\leq p\\\ell\ne j}}
       \frac{-q^\ell}{q^j-q^\ell}
 =\frac{(-1)^{p-j}q^{(p-j)(p-j+1)/2}}
       {(q;q)_j(q;q)_{p-j}}.
 \label{eq:qweights}
\end{equation}
These are simply the Lagrange interpolation weights for evaluation
at zero from the nodes $1,q,\ldots,q^p$.

\begin{theorem}[Exact reconstruction and uniform stability]\label{thm:richardson}
Let $x=a/2^d\in(-1,1)$, $p=\lfloor d/2\rfloor$, and $m\geq d+2$.
Then
\begin{equation}
 \boxed{\quad u(x)=\sum_{j=0}^p w_{p,j}(1/4)\,u_{m+j}(x).\quad}
 \label{eq:exact-richardson}
\end{equation}
At every order,
\begin{equation}
 \sum_{j=0}^p|w_{p,j}(1/4)|
 =\prod_{\ell=1}^p\frac{1+4^{-\ell}}{1-4^{-\ell}}<2.
 \label{eq:stable-weights}
\end{equation}
Thus absolute errors bounded by $\delta$ in each input spline value
produce an absolute output error smaller than $2\delta$.
\end{theorem}
\begin{proof}
Apply Lagrange interpolation to the polynomial
$z\mapsto R_x(4^{-m}z)$ from \cref{cor:poly-stabilize}. This proves
the equality.

Taking absolute values in \eqref{eq:qweights}, and using the finite
$q$-binomial product identity, gives
\begin{align}
 \sum_{j=0}^p|w_{p,j}(q)|
 &=\frac1{(q;q)_p}\sum_{k=0}^p
      q^{k(k+1)/2}\frac{(q;q)_p}{(q;q)_k(q;q)_{p-k}}\nonumber\\
 &=\frac{\prod_{\ell=1}^p(1+q^\ell)}{(q;q)_p}.
 \label{eq:qnorm-proof}
\end{align}
The finite product identity follows by expanding
$\prod_{\ell=1}^p(1+tq^\ell)$ and comparing the coefficient of
$t^k$; alternatively its coefficients satisfy the same two-term
recurrence on adjoining the last factor.

For the strict bound, $p=0,1$ are immediate. At $q=1/4$, the first
two factors multiply to $17/9$. For the remaining factors use
\begin{equation}
 \sum_{\ell=3}^{\infty}\log\frac{1+q^\ell}{1-q^\ell}
 \leq\frac{2q^3}{(1-q)(1-q^6)}=\frac{512}{12285}.
\end{equation}
Since $e^t\leq(1-t)^{-1}$ for $0\leq t<1$, the entire product is
at most
\begin{equation}
 \frac{17}{9}\frac1{1-512/12285}
 =\frac{208845}{105957}<2.
\end{equation}
Finally apply the triangle inequality to the input perturbations.
\end{proof}

The first three weight rows are
\begin{equation}
 (1),\qquad
 \left(-\frac13,\frac43\right),\qquad
 \left(\frac1{45},-\frac49,\frac{64}{45}\right).
\end{equation}
These familiar extrapolation weights are not new in isolation.
The point is the exact stopping theorem that makes them recover
physical-space dyadic values in finitely many levels.

\subsection{Examples and a closed finite formula}

\begin{example}
The following identities hold at every indicated sufficiently large level:
\begin{align}
 u_m(1/2)&=\frac12, &&m\geq3,\label{eq:examples-dyadic}\\
 u_m(1/4)&=\frac{67}{72}+\frac49\,4^{-m}, &&m\geq4,\nonumber\\
 u_m(1/8)&=\frac{287}{288}+\frac29\,4^{-m}, &&m\geq5,\nonumber\\
 u_m(1/16)&=\frac{2073457}{2073600}
          +\frac5{162}\,4^{-m}
          -\frac{3968}{2025}\,16^{-m}, &&m\geq6.\nonumber
\end{align}
For example, $u''(1/16)=-5/9$ and $u^{(4)}(1/16)=-1024$;
substitution into \eqref{eq:dyadic-backward} gives the last line.
The stated thresholds are uniform sufficient thresholds, not claims
of minimality for each example.
\end{example}

Combining \eqref{eq:exact-richardson} with
\eqref{eq:spline-positive} gives a closed finite expression for a
rational dyadic value. Put $m=d+2$ and $p=\lfloor d/2\rfloor$:
\begin{equation}
\begin{split}
 u(x)=\sum_{j=0}^p &\frac{w_{p,j}(1/4)\,2^{(m+j)(m+j-1)/2}}
                         {(m+j-1)!}\\[-2pt]
 &\times\sum_{n=0}^{2^{m+j}-1}\eps_n
 \left(x+1-2^{-(m+j)}-2^{1-m-j}n\right)_+^{m+j-1}.
\end{split}
\label{eq:closed-dyadic}
\end{equation}
Every quantity on the right is rational. This gives another proof
of classical dyadic rationality, and more importantly exhibits a
finite extrapolation mechanism with a uniform condition bound.
The formula is not advertised as the fastest existing algorithm.

\section{Repeated-prefix coefficients: an exact differential identity}\label{sec:lattice}

\subsection{The positive coefficient array and its natural normalization}

Let $c_m(k)$ be the coefficient of $z^k$ in $A_m(z)$ from
\eqref{eq:intro-Am}, extended by zero outside its support. Put
\begin{equation}
 D_m=2^m-m-1,\qquad B_m=A_m(1)=2^{m(m-1)/2}.
 \label{eq:DB}
\end{equation}
The support is $0\leq k\leq D_m$. The coefficients are positive
integers on this interval and are symmetric about $D_m/2$.

Repeated inclusive summation multiplies a generating function by
$(1-z)^{-1}$. Therefore these are precisely the $m$-fold cumulative
coefficients of the finite signed block, and
\begin{equation}
 c_m(k)=\sum_{n=0}^{\min(k,2^m-1)}
            \eps_n\binom{k-n+m-1}{m-1},\qquad k\geq0.
 \label{eq:c-binomial}
\end{equation}
For $k<2^m$ this agrees with the $m$-fold inclusive prefix of the
infinite Thue--Morse sequence. Beyond that range, the finite block
and infinite sequence should not be conflated.

Let $D_0^*,\ldots,D_{m-1}^*$ be independent random variables with
$D_j^*$ uniform on $\{0,\ldots,2^j-1\}$. Then
\begin{equation}
 \Prob\left(\sum_{j=0}^{m-1}D_j^*=k\right)=\frac{c_m(k)}{B_m}.
 \label{eq:c-probability}
\end{equation}
Set
\begin{equation}
 \Delta=2h=2^{1-m},\qquad
 x_{m,k}=\Delta(k-D_m/2)=-1+(m+1)h+2hk,
 \label{eq:lattice-nodes}
\end{equation}
and define the scaled lattice mass
\begin{equation}
 \rho_m(x_{m,k})=\frac{c_m(k)}{\Delta B_m}.
 \label{eq:rho}
\end{equation}
The division by $\Delta$ is essential: $\rho_m$ is a density-scale
quantity, not the probability of one lattice point.

\subsection{A centered factorial polynomial identity}

The next lemma is a useful form of the generalized Bernoulli
polynomial calculus; the standard generating function is recorded in
\cite{dlmf}. Its short proof is included to fix all normalizations.

\begin{lemma}[Centered factorial conversion]\label{lem:centered-factorial}
For $m\geq2$ define
\begin{equation}
 p_{m-1}(w)=\prod_{r=1}^{m-1}(w+r-m/2).
\end{equation}
As an identity of polynomials,
\begin{equation}
 p_{m-1}(w)=
 \left(\frac{\D_w/2}{\sinh(\D_w/2)}\right)^m w^{m-1}.
 \label{eq:centered-factorial}
\end{equation}
The power series on the right acts finitely: every derivative of
order greater than $m-1$ annihilates the polynomial.
\end{lemma}
\begin{proof}
By coefficient extraction and the formal change of variable
$y=e^z-1$,
\begin{align}
 [z^{m-1}]\left(\frac{z}{2\sinh(z/2)}\right)^m e^{wz}
 &=\Res_{z=0}\frac{e^{(w+m/2)z}}{(e^z-1)^m}\dd z\nonumber\\
 &=\Res_{y=0}\frac{(1+y)^{w+m/2-1}}{y^m}\dd y\nonumber\\
 &=\binom{w+m/2-1}{m-1}
 =\frac{p_{m-1}(w)}{(m-1)!}.
 \label{eq:residue-factorial}
\end{align}
The first coefficient equals $1/(m-1)!$ times the differential
operator in \eqref{eq:centered-factorial} applied to $w^{m-1}$.
The change of variable is valid formally because $e^z-1=z+O(z^2)$.
\end{proof}

\begin{theorem}[Exact lattice-to-spline differential conversion]\label{thm:lattice-operator}
For $m\geq2$ and every $k\in\Z$,
\begin{equation}
 \boxed{\quad
 \rho_m(x_{m,k})=
 \sum_{j=0}^{\lfloor(m-1)/2\rfloor}
        \gamma_j(m)h^{2j}u_m^{(2j)}(x_{m,k}).\quad}
 \label{eq:lattice-exact-operator}
\end{equation}
When $m$ is odd the highest derivative is a cellwise derivative,
but the lattice nodes are then cell centers. When $m$ is even the
highest derivative used is of order $m-2$ and is continuous.
Thus every displayed evaluation is unambiguous.
\end{theorem}
\begin{proof}
At $x=x_{m,k}$ put
\begin{equation}
 t=\frac{x+1-h}{2h}=k+\frac m2,
 \qquad J_m=2^{-(m-1)(m-2)/2}.
\end{equation}
The signed spline formula becomes
\begin{equation}
 u_m(x)=\frac{J_m}{(m-1)!}
       \sum_{\substack{0\leq n<2^m\\n<t}}\eps_n(t-n)^{m-1}.
 \label{eq:spline-rescaled-t}
\end{equation}
On the other hand, \eqref{eq:c-binomial} and
$1/(\Delta B_m)=J_m$ give
\begin{equation}
 \rho_m(x)=\frac{J_m}{(m-1)!}
       \sum_{\substack{0\leq n<2^m\\n\leq k}}\eps_n p_{m-1}(t-n).
 \label{eq:rho-rescaled-t}
\end{equation}
One can extend this last sum from $n\leq k$ to $n<t$: every new
integer has $n=k+i$ with $1\leq i<m/2$, and then
$p_{m-1}(m/2-i)=0$. The same observation works if the original
sum is empty or if the finite block truncates the range.

Now apply \cref{lem:centered-factorial} term by term to
\eqref{eq:spline-rescaled-t}. Since $\D_t=2h\D_x$, its operator
becomes $(h\D_x/\sinh(h\D_x))^m$, whose coefficients are
$\gamma_j(m)h^{2j}$. This proves the formula. At an even-level
knot, the derivatives involved have positive remaining polynomial
degree; the limiting formula agrees from both sides.
\end{proof}

The first terms are
\begin{equation}
 \rho_m(x)=u_m(x)-\frac m6h^2u_m''(x)
        +\left(\frac{m^2}{72}+\frac m{180}\right)h^4u_m^{(4)}(x)
        +\cdots,
 \label{eq:operator-first}
\end{equation}
where the complete expression is the finite sum in
\eqref{eq:lattice-exact-operator}. This identity is stronger than a
formal differential approximation.

\section{The all-orders uniform local-limit theorem}\label{sec:local-limit}

The exact operator in \cref{thm:lattice-operator} is useful
algebraically. To obtain a uniform fixed-order remainder as $m$ grows,
it is cleaner to work on the lattice Fourier interval.

\subsection{Exact removal of the lattice}

\begin{proposition}[Continuous noise representation]\label{prop:dequantization}
Let
\begin{equation}
 X_m=\Delta\left(\sum_{j=0}^{m-1}D_j^*-D_m/2\right).
\end{equation}
There is a random variable $W_m$, independent of $X_m$, such that
\begin{equation}
 Y\ \overset{\mathrm d}=\ X_m+hW_m,
 \qquad W_m=V_1+\cdots+V_m+Y',
 \label{eq:dequantization}
\end{equation}
where all summands on the right are independent, the $V_j$ are
uniform on $[-1,1]$, and $Y'$ is a copy of $Y$. In particular
\begin{equation}
 \operatorname{Var}(W_m)=\frac m3+\frac19.
 \label{eq:noise-variance}
\end{equation}
The characteristic function of $X_m$ satisfies
\begin{equation}
 \psi_m(t)=\frac{Q(t)}{Q(ht)\sinc(ht)^m}
 \qquad\left(|t|\leq\frac{\pi}{2h}\right).
 \label{eq:psi-quotient}
\end{equation}
\end{proposition}
\begin{proof}
Adjoin independent $U_j$ uniform on $[0,1]$ to $D_j^*$. Then
$D_j^*+U_j$ is uniform on $[0,2^j]$, and
\begin{equation}
 \Delta\left(D_j^*+U_j-2^{j-1}\right)
 =\Delta\left(D_j^*-\frac{2^j-1}{2}\right)+h(2U_j-1).
\end{equation}
Summing over $j$ produces $X_m$ plus $m$ independent uniform
noises of half-width $h$. The continuous sum has the law of $Y_m$
after reversing the order of its widths. Adjoin the tail $hY'$.
The variance follows from $\mu_2=1/9$. Taking characteristic
functions proves the quotient wherever its denominator is nonzero,
in particular on the stated interval.
\end{proof}

This calculation explains the leading factor $m$. Continuous spline
truncation omits one self-similar tail of scale $h$, whereas converting
a lattice sum to its continuous counterpart requires $m$ additional
small uniforms. Their combined variance is of order $mh^2$.

\subsection{Uniform expansion and an explicit remainder constant}

For use in the proof, write
\begin{equation}
 \mathcal H_m(z)=\sum_{j\geq0}H_j(m)z^j
 =\exp\left(\sum_{r\geq1}\lambda_r
               \left(m+\frac1{4^r-1}\right)z^r\right).
\end{equation}
By \eqref{eq:psi-quotient},
$\psi_m(t)=Q(t)\mathcal H_m(h^2t^2)$ on the fundamental interval.

\begin{theorem}[Uniform all-orders local-limit expansion]\label{thm:local-limit}
For each fixed integer $N\geq0$, there exists an explicit constant
$C_N$ such that for every $m\geq4$,
\begin{equation}
 \boxed{
 \sup_{k\in\Z}\left|
 \rho_m(x_{m,k})-
 \sum_{j=0}^N(-1)^jH_j(m)h^{2j}u^{(2j)}(x_{m,k})
 \right|
 \leq C_N(m+1)^{N+1}h^{2N+2}.}
 \label{eq:uniform-local-limit}
\end{equation}
One permissible, deliberately nonoptimized choice is obtained by
putting $L=4N+8$ and taking
\begin{equation}
 C_N=\frac2\pi\left\{
 \frac1{2N+3}+\frac{K_L}{2N+5}
 +\frac{K_L}{4N+7}
 +K_L\sum_{j=0}^N\frac1{4N+7-2j}
 \right\}.
 \label{eq:CN-explicit}
\end{equation}
\end{theorem}
\begin{proof}
Lattice Fourier inversion and our normalization give
\begin{equation}
 \rho_m(x_{m,k})=\frac1{2\pi}
 \int_{-\pi/(2h)}^{\pi/(2h)}e^{-\ii t x_{m,k}}\psi_m(t)\dd t.
 \label{eq:lattice-inversion}
\end{equation}
The finite proposed approximation is the full-line inverse Fourier
transform of
$Q(t)\sum_{j=0}^NH_j(m)h^{2j}t^{2j}$, because multiplication by
$t^{2j}$ corresponds to $(-1)^j\D^{2j}$.
We estimate the difference of these transforms in $L^1$; the resulting
bound is independent of $k$.

\emph{Low frequencies.}
Put $T=h^{-1/2}$ and $R=(m+1)^{-1}$. Since
$\lambda_r\leq2/(r\pi^{2r})$, positivity gives
\begin{equation}
 \log\mathcal H_m(R)
 \leq\frac2{\pi^2-R}\leq\frac2{\pi^2-1}<\log2.
\end{equation}
Hence
\begin{equation}
 \mathcal H_m(R)<2,
 \qquad H_j(m)\leq2(m+1)^j.
 \label{eq:H-coeff-bound}
\end{equation}
For $|t|\leq T$, put $z=h^2t^2\leq h$. Since $m\geq4$,
$z/R\leq(m+1)h\leq5/16<1$. Positivity of all coefficients implies
\begin{equation}
 \left|\mathcal H_m(z)-\sum_{j=0}^NH_j(m)z^j\right|
 \leq2\bigl((m+1)z\bigr)^{N+1}.
 \label{eq:H-remainder}
\end{equation}
Therefore the low-frequency error before division by $2\pi$ is at most
\begin{equation}
 2(m+1)^{N+1}h^{2N+2}
       \int_{\R}|t|^{2N+2}|Q(t)|\dd t.
 \label{eq:low-error}
\end{equation}
Using \cref{lem:decay} with $L=4N+8$ bounds the integral by
\begin{equation}
 2\left(\frac1{2N+3}+\frac{K_L}{2N+5}\right).
 \label{eq:weighted-integral-bound}
\end{equation}

\emph{High frequencies of the exact lattice transform.}
For $|s|\leq\pi/2$, $\sinc s\geq2/\pi$. Also, multiplying the
bounds $\sinc(2^{-j}s)\geq1-2^{-2j}s^2/6$ gives
\begin{equation}
 Q(s)\geq1-\frac{s^2}{18}\geq1-\frac{\pi^2}{72}>\frac12.
\end{equation}
Thus on the fundamental interval,
\begin{equation}
 |\psi_m(t)|\leq2(\pi/2)^m|Q(t)|.
\end{equation}
Write $\beta=\log_2(\pi/2)<1$. Integration over $|t|>T$ yields
\begin{equation}
 \int_{T<|t|\leq\pi/(2h)}|\psi_m(t)|\dd t
 \leq\frac{4K_L}{L-1}h^{(L-1)/2-\beta}
 \leq\frac{4K_L}{4N+7}h^{2N+2}.
 \label{eq:psi-tail-bound}
\end{equation}

\emph{High frequencies of the finite approximation.}
By \eqref{eq:H-coeff-bound} and \cref{lem:decay},
\begin{align}
 &\int_{|t|>T}|Q(t)|\sum_{j=0}^NH_j(m)h^{2j}|t|^{2j}\dd t\nonumber\\
 &\quad\leq4K_L\sum_{j=0}^N
       \frac{(m+1)^jh^{(L-1)/2+j}}{L-2j-1}\nonumber\\
 &\quad\leq4K_L(m+1)^{N+1}h^{2N+2}
          \sum_{j=0}^N\frac1{4N+7-2j}.
 \label{eq:approx-tail-bound}
\end{align}
Add \eqref{eq:low-error}--\eqref{eq:approx-tail-bound}, divide by
$2\pi$, and use $(m+1)^{N+1}\geq1$ in the exact-transform tail.
The result is precisely \eqref{eq:uniform-local-limit} with
\eqref{eq:CN-explicit}.
\end{proof}

\begin{remark}[Meaning of the local limit]
The supremum in \eqref{eq:uniform-local-limit} includes nodes outside
the finite support, where $c_m(k)=0$. This is a uniform comparison
of scaled point masses with a density evaluated at the same nodes.
It is not convergence in total variation of a discrete measure to
an absolutely continuous one.
\end{remark}

\begin{corollary}[Sharp first local-limit error]\label{cor:local-sharp}
Uniformly for $k\in\Z$,
\begin{equation}
 \frac{\rho_m(x_{m,k})-u(x_{m,k})}{m4^{-m}}
       +\frac16u''(x_{m,k})\longrightarrow0.
 \label{eq:normalized-lattice-error}
\end{equation}
In particular,
\begin{equation}
 \lim_{m\to\infty}
 \frac{\sup_{k\in\Z}|\rho_m(x_{m,k})-u(x_{m,k})|}{m4^{-m}}
 =\frac43.
 \label{eq:lattice-sharp-constant}
\end{equation}
\end{corollary}
\begin{proof}
Use \cref{thm:local-limit} with $N=1$ and
$H_1(m)/m\to1/6$. Its remainder divided by $m4^{-m}$ tends to
zero. The node spacing tends to zero, so the sampled supremum of
the continuous compactly supported function $|u''|$ tends to
$\norm{u''}_\infty=8$ by \cref{prop:copies}.
\end{proof}

\section{Exact dyadic stopping for the integer coefficient array}\label{sec:dyadic-lattice}

The local-limit expansion has an unexpectedly exact specialization.
It is obtained from the finite operator, not by declaring an
asymptotic series to terminate.

\begin{theorem}[Exact dyadic lattice expansion]\label{thm:dyadic-lattice}
Fix $x=a/2^d\in(-1,1)$. If $m$ is odd and $m\geq d+2$, then
\begin{equation}
 k_m=2^{m-1}(1+x)-\frac{m+1}{2}\in\Z,
 \qquad x_{m,k_m}=x,
\end{equation}
and
\begin{equation}
 \boxed{\quad
 \rho_m(x)=\sum_{j=0}^{\lfloor d/2\rfloor}
                 (-1)^jH_j(m)4^{-jm}u^{(2j)}(x).
 \quad}
 \label{eq:exact-dyadic-lattice}
\end{equation}
The right-hand side is an exact finite sum of polynomials in $m$
times powers of $4^{-m}$.
\end{theorem}
\begin{proof}
The stated parity makes $k_m$ an integer. On the cell about $x$,
$u_m$ has degree at most $d$ by \cref{lem:dyadic-cell}. Hence the
finite differential conversion in \cref{thm:lattice-operator} has
no nonzero terms with $2j>d$ at this point. Substitute
\eqref{eq:dyadic-backward} into its remaining terms. The coefficient
of $h^{2j}u^{(2j)}(x)$ is the convolution
$\sum_{\ell=0}^j\gamma_\ell(m)b_{j-\ell}$, which is
$(-1)^jH_j(m)$ by \eqref{eq:coefficient-convolution}.
\end{proof}

\begin{example}[A closed repeated-prefix identity]\label{ex:quarter-lattice}
For every odd $m\geq5$,
\begin{equation}
 \rho_m(1/4)=\frac{67}{72}
       +\left(\frac{4m}{3}+\frac49\right)4^{-m}.
 \label{eq:quarter-lattice-exact}
\end{equation}
Equivalently, an explicitly indexed coefficient of the positive
Thue--Morse quotient is
\begin{equation}
 \boxed{\quad
 c_m\left(5\cdot2^{m-3}-\frac{m+1}{2}\right)
 =2^{m(m-1)/2-m+1}
  \left\{\frac{67}{72}+
         \left(\frac{4m}{3}+\frac49\right)4^{-m}\right\}.
 \quad}
 \label{eq:closed-prefix-quarter}
\end{equation}
This also evaluates the corresponding $m$-fold inclusive prefix
of the infinite sequence, since the index is below $2^m$.
\end{example}
\begin{proof}
At $x=1/4$, $u=67/72$, $u''=-8$, and all derivatives of order
above two vanish. Apply \cref{thm:dyadic-lattice} and then undo
the density normalization.
\end{proof}

At $x=1/16$ the exact identity has three terms:
\begin{equation}
 \rho_m(1/16)=\frac{2073457}{2073600}
       +\frac59H_1(m)4^{-m}
       -1024H_2(m)16^{-m},
 \qquad m\geq7\text{ odd}.
 \label{eq:sixteenth-lattice}
\end{equation}
The increasing polynomial factors $H_j(m)$ are therefore not
merely artifacts of a Fourier proof; they occur in exact rational
coefficient identities.

\paragraph{A useful distinction.}
For the spline values at a fixed dyadic point, the coefficients of
$4^{-jm}$ are constants depending on that point. For the discrete
array values, the coefficients are polynomials in $m$. The extra
polynomials come from the $m$ independent lattice-removal noises.
This distinction prevents the incorrect use of the simple spline
Richardson weights as an exact annihilator of all lattice errors.

\section{Fourier-energy moments and lognormal twins}\label{sec:spectral}

We now use the $L^2$ version of the derivative-copy identity. This
has a different flavor from either approximation theorem.

\begin{definition}[The normalized energy measure]
Let
\begin{equation}
 I_0=\int_{\R}|Q(t)|^2\dd t=2\pi\norm{u}_2^2,
 \qquad \dd\nu(t)=I_0^{-1}|Q(t)|^2\dd t.
 \label{eq:nu-def}
\end{equation}
This is an absolutely continuous probability measure on the angular
frequency line. It is not the singular diffraction measure of the
unscaled Thue--Morse substitution system.
\end{definition}

\begin{theorem}[Exact energy moments]\label{thm:energy-moments}
Every moment of $\nu$ exists, and for $r\geq0$,
\begin{equation}
 \int_{\R}t^{2r}\dd\nu(t)=2^{r(r+1)},\qquad
 \int_{\R}t^{2r+1}\dd\nu(t)=0.
 \label{eq:energy-moments}
\end{equation}
\end{theorem}
\begin{proof}
The decay estimate \eqref{eq:Qbound} gives existence of every
moment. Evenness gives the odd moments. Plancherel's identity and
\eqref{eq:Lp-copies} at $p=2$ give
\begin{equation}
 \frac{\int t^{2r}|Q(t)|^2\dd t}{\int|Q(t)|^2\dd t}
 =\frac{\norm{u^{(r)}}_2^2}{\norm{u}_2^2}
 =K_r^2=2^{r(r+1)}.
\end{equation}
\end{proof}

\subsection{An explicit absolutely continuous twin}

\begin{theorem}[Lognormal moment twin]\label{thm:lognormal}
Let $S$ take values $-1,1$ with equal probabilities and let
$L\sim N((\log2)/2,(\log2)/2)$ be independent of $S$. Then the
law of $T=S e^L$ has the same moments as $\nu$, but it is a
different probability measure. Its density is
\begin{equation}
 g(t)=\frac1{2|t|\sqrt{\pi\log2}}
       \exp\left(-\frac{(\log|t|-(\log2)/2)^2}{\log2}\right),
 \quad t\ne0,
 \qquad g(0)=0.
 \label{eq:lognormal-density}
\end{equation}
Consequently the Hamburger moment problem for
\eqref{eq:energy-moments} is indeterminate.
\end{theorem}
\begin{proof}
For a normal random variable of mean $\mu$ and variance $\sigma^2$,
completion of the square in its density gives
$\E e^{sL}=e^{s\mu+s^2\sigma^2/2}$. At $s=2r$ and the stated
parameters, this is $2^{r(r+1)}$. The independent sign kills all
odd moments. Change of variable $t=e^L$, followed by symmetrization,
gives \eqref{eq:lognormal-density}.

The density $g$ extends continuously to zero with value zero,
whereas the density of $\nu$ has value $1/I_0>0$ there. They
therefore differ on a neighborhood of zero, not merely at a single
point. Thus two distinct measures have the same complete moment
sequence.
\end{proof}

Lognormal moment indeterminacy is classical \cite{heyde}. The
contribution here is the exact identification of a compactly
supported smooth function's Fourier-energy measure with that
moment class. No inference from the failure of a sufficient
uniqueness criterion is needed: two distinct measures have been
constructed explicitly.

\subsection{A continuum of atomic twins}

For $\alpha\in[0,1)$ set
\begin{equation}
 \Theta(\alpha)=\sum_{k\in\Z}2^{-(k+\alpha)^2}.
 \label{eq:theta}
\end{equation}
It is finite and strictly positive.

\begin{theorem}[Theta-supported moment twins]\label{thm:theta}
Every measure
\begin{equation}
 \nu_\alpha=\frac1{2\Theta(\alpha)}
 \sum_{k\in\Z}2^{-(k+\alpha)^2}
 \left(\delta_{2^{k+\alpha+1/2}}+
       \delta_{-2^{k+\alpha+1/2}}\right)
 \label{eq:atomic-twin}
\end{equation}
is a probability measure with the moment sequence
\eqref{eq:energy-moments}.
\end{theorem}
\begin{proof}
The total mass is one and the odd moments vanish. For the even
moments, complete the square in the exponent:
\begin{align}
 \int t^{2r}\dd\nu_\alpha(t)
 &=\frac1{\Theta(\alpha)}
     \sum_{k\in\Z}2^{-(k+\alpha)^2+2r(k+\alpha+1/2)}\nonumber\\
 &=\frac{2^{r^2+r}}{\Theta(\alpha)}
     \sum_{k\in\Z}2^{-(k+\alpha-r)^2}
 =2^{r(r+1)}.
\end{align}
The last equality shifts the integer summation index by $r$.
All sums converge absolutely.
\end{proof}

The theta construction is a standard type of discrete lognormal
moment mechanism; the displayed proof is included rather than
claiming that mechanism itself as a new discovery. Its combination
with the energy identity gives the exact polynomial quadrature
\begin{equation}
 \frac1{I_0}\int_{\R}p(t)|Q(t)|^2\dd t
 =\int_{\R}p(t)g(t)\dd t
 =\int_{\R}p(t)\dd\nu_\alpha(t)
 \quad\text{for every polynomial }p.
 \label{eq:polynomial-quadrature}
\end{equation}
This equality cannot be extended to all bounded continuous test
functions, since the measures are distinct.

\subsection{Identical jets, and what restores uniqueness}

Define the normalized autocorrelation
\begin{equation}
 R_u(x)=\frac{\int_{\R}u(y)u(y+x)\dd y}{\norm{u}_2^2}.
 \label{eq:autocorrelation}
\end{equation}
It is a $C^\infty$ function supported on $[-2,2]$ and is the
characteristic function of $\nu$. The energy moments give
\begin{equation}
 R_u^{(2r)}(0)=(-1)^r2^{r(r+1)},\qquad R_u^{(2r+1)}(0)=0.
 \label{eq:autocorrelation-jet}
\end{equation}
Its formal Taylor series at zero has radius of convergence zero,
since $2^{r(r+1)}/(2r)!$ has unbounded $2r$th root. The
characteristic functions of the lognormal and atomic twins have
exactly the same derivatives at zero. Thus their differences from
$R_u$ are smooth and flat at zero, but not identically zero.

Moments alone lose information, but the original scaling relation
retains it.
\begin{proposition}[Refinement restores uniqueness]\label{prop:refinement-unique}
Suppose a probability density $f$ is continuous at zero and satisfies
pointwise
\begin{equation}
 f(2t)=\sinc^2(t)f(t),\qquad t\in\R.
 \label{eq:energy-refinement}
\end{equation}
Then $f(t)=I_0^{-1}Q(t)^2$.
\end{proposition}
\begin{proof}
Iterating toward zero gives
\begin{equation}
 f(t)=f(2^{-n}t)\prod_{j=1}^n\sinc^2(2^{-j}t).
\end{equation}
Continuity at zero and convergence of the product imply
$f(t)=f(0)Q(t)^2$. Integration fixes $f(0)=1/I_0$.
\end{proof}

\section{A separated geometric extension}\label{sec:geometric}

The moment phenomenon is not an accidental equality special to one
choice of units. It follows from separation of derivative copies.
Here is a one-parameter extension that makes this explicit.

For $0<\theta\leq1/2$, let $v_\theta$ be the density of
\begin{equation}
 Y_\theta=(1-\theta)\sum_{j=0}^{\infty}\theta^j V_j.
 \label{eq:geometric-series}
\end{equation}
It is supported on $[-1,1]$ and is smooth by the same Fourier
argument as before. At $\theta=1/2$ it equals $u$.

\begin{theorem}[Separated-copy norms and their moment class]\label{thm:geometric}
For $r\geq0$ and $1\leq p\leq\infty$,
\begin{equation}
 \norm{v_\theta^{(r)}}_p
 =(2\theta)^{r/p}[2(1-\theta)]^{-r}
       \theta^{-r(r+1)/2}\norm{v_\theta}_p,
 \label{eq:geometric-norm}
\end{equation}
where $(2\theta)^{r/p}=1$ at $p=\infty$.
If $\nu_\theta$ is its normalized Fourier-energy measure, then
\begin{equation}
 \int t^{2r}\dd\nu_\theta(t)
 =[2(1-\theta)^2]^{-r}\theta^{-r^2},
 \qquad \int t^{2r+1}\dd\nu_\theta(t)=0.
 \label{eq:geometric-moments}
\end{equation}
These are the moments of a symmetric lognormal variable with
logarithmic mean and variance
\begin{equation}
 \mu_\theta=-\frac12\log\bigl(2(1-\theta)^2\bigr),
 \qquad \sigma_\theta^2=-\frac12\log\theta.
 \label{eq:geometric-lognormal}
\end{equation}
\end{theorem}
\begin{proof}
Differentiating the first uniform factor yields
\begin{equation}
 v_\theta'(x)=\frac1{2(1-\theta)\theta}
 \left\{v_\theta\left(\frac{x+1-\theta}{\theta}\right)
       -v_\theta\left(\frac{x-1+\theta}{\theta}\right)\right\}.
 \label{eq:geometric-refinement}
\end{equation}
The two copies have disjoint interiors because $\theta\leq1/2$.
At depth $r$ there are $2^r$ disjoint copies, each of horizontal
scale $\theta^r$, with signs given by binary endpoint parity.
Their common vertical multiplier is
\begin{equation}
 [2(1-\theta)]^{-r}\theta^{-1-2-\cdots-r}.
\end{equation}
Integrating the $p$th power gives \eqref{eq:geometric-norm}.
At $p=2$, squaring the norm ratio simplifies to
$[2(1-\theta)^2]^{-r}\theta^{-r^2}$. Plancherel proves the
moment formula. The normal exponential-moment calculation gives
\eqref{eq:geometric-lognormal}.
\end{proof}

For a completely explicit atomic version, put $a=-\log\theta>0$,
$\mu=\mu_\theta$, and
\begin{equation}
 \Theta_a(\alpha)=\sum_{k\in\Z}e^{-a(k+\alpha)^2}.
\end{equation}
Then
\begin{equation}
 \frac1{2\Theta_a(\alpha)}\sum_{k\in\Z}e^{-a(k+\alpha)^2}
 \left(\delta_{e^{\mu+a(k+\alpha)}}+
       \delta_{-e^{\mu+a(k+\alpha)}}\right)
 \label{eq:general-atomic}
\end{equation}
has the same moments, by shifting $k$ after completing the square.
This generalization concerns separated copies. For $\theta>1/2$
the copies overlap, the simple norm calculation no longer applies,
and \eqref{eq:geometric-moments} is not asserted.

\section{Exact computations and reproducibility}\label{sec:computations}

The accompanying \texttt{verify.py} uses only the Python standard
library. Algebraic tests use arbitrary-size integers and
\texttt{fractions.Fraction}, not floating-point agreement. The saved
\texttt{verification\_results.json} records the actual run used in
preparing this article.

\subsection{What was checked}

The script performed 3,796 categorized exact checks: 35 Prouhet
moment identities, 135 derivative-copy identities, 2,307 dyadic
polynomial identities, 127 exact Richardson reconstructions, 510
prefix-convolution identities, 15 extrapolation-norm identities,
501 lattice differential conversions, and 166 dyadic lattice
identities. The tested ranges and loop bounds are visible in the
script. Some checks compare formulas obtained from the same
underlying theory; they are useful for catching sign, indexing,
and normalization errors, not substitutes for the general proofs.

The implementation constructs the positive coefficient array by
sliding-window convolution with all-ones vectors. At level $j$ the
vector length is $O(2^j)$, so all levels through $m$ use $O(2^m)$
integer operations. This is an arithmetic-operation count, not a
bit-complexity bound: the integers themselves grow.

\subsection{A nondyadic corrected approximation}

For $x=1/3$, the reference value was computed from $C_{18,6}(x)$.
The exact global certificate from \cref{thm:global} is less than
$3.469\times10^{-59}$, and the reference starts
\begin{equation}
 u(1/3)=0.8198348851985181\ldots.
\end{equation}
The rational reference and its rational error bound are included
in the JSON output. The entries in \cref{tab:corrections} are
rounded absolute differences from that reference; its certified
uncertainty is negligible at the digits displayed.

\begin{table}[htbp]
\centering
\renewcommand{\arraystretch}{1.2}
\begin{tabular}{rrrr}
\toprule
$m$ & $|u-u_m|$ & $|u-C_{m,1}|$ & $|u-C_{m,2}|$\\
\midrule
6 & $8.90547\cdot10^{-5}$ & $5.96017\cdot10^{-8}$ & ---\\
8 & $5.56024\cdot10^{-6}$ & $2.29567\cdot10^{-10}$ & $8.44219\cdot10^{-14}$\\
10 & $3.47493\cdot10^{-7}$ & $8.95576\cdot10^{-13}$ & $2.03133\cdot10^{-17}$\\
12 & $2.17182\cdot10^{-8}$ & $3.49806\cdot10^{-15}$ & $4.95223\cdot10^{-21}$\\
\bottomrule
\end{tabular}
\caption{Corrected spline errors at $x=1/3$. The omitted entry has
not been used because the uniform classical-regularity threshold
for $N=2$ is $m\geq8$.}\label{tab:corrections}
\end{table}

For example, at $m=12$, $N=2$, the actual displayed error is
approximately $4.95223\times10^{-21}$, while the rigorous global
bound is approximately $6.03996\times10^{-21}$. The same certificate
works at every real $x$, not only at this sample point.

\subsection{An exact discrete-array experiment}

At $x=1/4$ and odd levels the normalized error is exactly
$4/3+4/(9m)$ by \eqref{eq:quarter-lattice-exact}. The script
computed the coefficient arrays independently by positive
convolution and obtained \cref{tab:lattice}. In every row the
residual after subtracting the two terms of
\eqref{eq:quarter-lattice-exact} is exactly zero.

\begin{table}[htbp]
\centering
\renewcommand{\arraystretch}{1.25}
\begin{tabular}{rrr}
\toprule
$m$ & $\rho_m(1/4)$ & $\dfrac{\rho_m(1/4)-67/72}{m4^{-m}}$\\[4pt]
\midrule
5 & $15/16$ & $64/45$\\
7 & $1907/2048$ & $88/63$\\
9 & $15247/16384$ & $112/81$\\
11 & $487881/524288$ & $136/99$\\
13 & $1951517/2097152$ & $160/117$\\
15 & $124897055/134217728$ & $184/135$\\
\bottomrule
\end{tabular}
\caption{Exact repeated-prefix coefficients after density
normalization. The last column tends to $4/3$.}\label{tab:lattice}
\end{table}

The script also makes finite numerical checks of the theta moment
formulas at several values of $\alpha$. Those checks are explicitly
numerical; the completion-of-the-square proof, not truncation of a
theta sum, establishes the exact identity for every moment.

\section{Consequences, limitations, and concrete next problems}\label{sec:frontiers}

The strongest conclusion for the motivating approximation question
is that fixed-order strong spline corrections need no deleted knot
neighborhoods. The sufficient condition is a regularity threshold
on the spline level. This directly addresses the strong-spline
conjecture in the supplied atlas \cite{atlas}; it does not establish
that document's separate relative-error endpoint saddle evaluator.

The exact dyadic results exhibit a second phenomenon. A function
can be globally nonanalytic and yet admit finite exact extraction
from an approximation sequence at every dyadic point. The source
of this exactness is a low-degree \emph{finite spline cell}, not
analyticity of the limit. The highest nonzero derivative even has
an explicit Thue--Morse sign in
\eqref{eq:highest-dyadic-derivative}.

On the discrete side, the exact centered-factorial identity
separates two sources of correction: $\gamma_j(m)$ accounts for
lattice removal, and $b_j$ accounts for the infinite continuous
tail. Their convolution produces $H_j(m)$. This is an effective
way to organize further coefficient identities without repeatedly
rediscovering their signs and centering terms.

Finally, the spectral calculation shows that an entire family of
energy moments can conceal major differences between probability
measures. A continuous density positive at zero, a lognormal-type
density flat at zero, and a purely atomic measure can share every
polynomial moment. The refinement equation restores information
that the moments alone discard.

\begin{problem}[Relative endpoint accuracy]
Develop bounds for $C_{m,N}(x)-u(x)$ relative to $u(x)$ as $x$
approaches an endpoint while $m$ and $N$ vary. The present global
absolute bounds remain true there, but division by the very small
value of $u(x)$ can make them ineffective. A proof will need local
information beyond a global derivative norm.
\end{problem}

\begin{problem}[Efficient exact extraction]
Determine the best bit complexity of dyadic reconstruction from
positive coefficient arrays, local spline cells, or the signed
formula \eqref{eq:closed-dyadic}. The extrapolation condition bound
is uniform, but the cost of producing the inputs and the cost of
large integer cancellation must be analyzed separately.
\end{problem}

\begin{problem}[Overlapping geometric copies]
For $1/2<\theta<1$, compute or characterize the cross terms in
$\norm{v_\theta^{(r)}}_2^2$. The separated calculation in
\cref{thm:geometric} identifies exactly where the argument changes:
the derivative-copy intervals now overlap. It is natural to ask
which parameters, if any, retain an explicit moment class.
\end{problem}

\subsection{A formalization route without a formalization claim}

The main dependency chain is short. Finite convolution and binary
endpoint enumeration yield the derivative-copy identity. Disjoint
supports give its norm formula. A Taylor remainder bound and the
probability-tail identity then yield the global theorem. The
dyadic cell lemma and finite triangular inversion are algebraic
once that bridge is available. Separately, the centered-factorial
identity can be formalized in a polynomial ring, and the lattice
local-limit theorem requires Fourier inversion with explicit
integrable majorants. Plancherel plus disjoint supports proves
the energy moments.

These are proposed proof components, not assertions that new Lean
theorems have been compiled or added to the repository. The
manuscript and the exact-arithmetic script are the deliverables
of this investigation.

\appendix
\section{Normalization dictionary and implementation details}\label{app:normalization}

\subsection{Frequency conventions}

The angular-frequency transform used here is $Q(t)$ in
\eqref{eq:Q}. Under the cyclic-frequency convention
$\widehat u(\xi)=\int e^{-2\pi\ii\xi x}u(x)\dd x$, evenness gives
\begin{equation}
 \widehat u(\xi)=Q(2\pi\xi)
   =\prod_{j=0}^{\infty}\sinc(\pi\xi/2^j).
 \label{eq:cyclic-frequency}
\end{equation}
A finite cyclic product indexed $j=0,\ldots,n$ therefore corresponds
to the spline $u_{n+1}$ in this article. Its omitted-tail scale is
$2^{-(n+1)}$, not $2^{-n}$. The even moments of its normalized
cyclic-frequency energy measure are
\begin{equation}
 (2\pi)^{-2r}2^{r(r+1)}.
\end{equation}
This dictionary is important when comparing formulas with
repository documents that use a different Fourier convention or
start their product index at zero.

\subsection{Integer-only evaluation of a spline derivative}

For $x=a/b$ with $b>0$ and $p=m-1-r\geq1$, put
\begin{equation}
 T=a2^m+b(2^m-1).
\end{equation}
Then \eqref{eq:spline-deriv-positive} is
\begin{equation}
 u_m^{(r)}(a/b)
 =\frac{2^{m(m-1)/2}}{p!(b2^m)^p}
       \sum_{\substack{0\leq n<2^m\\T-2bn>0}}
           \eps_n(T-2bn)^p.
 \label{eq:integer-spline-evaluation}
\end{equation}
The sum is an integer and can be accumulated before one final
rational division. For a highest cellwise derivative, $p=0$,
use the same formula away from knots with the positive-part
indicator convention. The supplied implementation treats values
outside the support as zero before forming the signed sum.

\subsection{Reproducing the computation and PDF}

From the directory containing the archive contents, run
\begin{lstlisting}
python verify.py
pdflatex Thue_Morse_New_Deductions.tex
pdflatex Thue_Morse_New_Deductions.tex
pdflatex Thue_Morse_New_Deductions.tex
\end{lstlisting}
The Python script requires Python~3.10 or later and no third-party
packages. The LaTeX source uses standard mathematical packages and
the New~TX fonts available in a full TeX~Live installation. The
article source is self-contained: it needs no external figures,
bibliography database, generated tables, or source-code inclusion
to compile.

\section{Statement ledger}\label{app:ledger}

\begin{longtable}{@{}p{0.26\textwidth}p{0.67\textwidth}@{}}
\toprule
Result & Precisely what is established\\
\midrule
\endhead
Global correction &
\Cref{thm:global,thm:sharp,cor:holder}: all fixed orders, global
classical and H\"older norms, explicit constants, and sharp leading
supremum asymptotics.\\[4pt]
Increasing order &
\Cref{cor:growing}: a supergeometric absolute error certificate
with a specified admissible order; no relative endpoint or
numerical-conditioning claim.\\[4pt]
Dyadic splines &
\Cref{thm:dyadic,cor:exact-degree,thm:richardson}: exact polynomial
stabilization, exact degree in lowest terms, finite extraction,
and an order-independent absolute amplification bound.\\[4pt]
Discrete array &
\Cref{thm:lattice-operator,thm:local-limit,thm:dyadic-lattice}:
finite differential conversion, a uniform all-orders local-limit
bound, and exact stopping on fixed dyadic odd-level subsequences.\\[4pt]
Energy moments &
\Cref{thm:energy-moments,thm:lognormal,thm:theta}: exact moments,
explicit distinct continuous and atomic moment twins, and hence
moment indeterminacy.\\[4pt]
Geometric extension &
\Cref{thm:geometric}: exact separated-copy norms and the resulting
moment class for $0<\theta\leq1/2$, not for overlapping copies.\\[4pt]
Priority and verification &
Proofs and finite exact computations are supplied; universal
publication priority, peer review, and Lean formalization are
not claimed.\\
\bottomrule
\end{longtable}

\clearpage
\begin{thebibliography}{9}
\addcontentsline{toc}{section}{References}

\bibitem{atlas}
ProveIt project.
\emph{The Thue--Morse Sequence: Formula Atlas and Fabius--Rvachev
Frontier Results}.
Research draft, consolidated 28 August 2026, in the repository
maintained by Vladimir Reshetnikov. Source inspected 4 September 2026.
\url{https://github.com/VladimirReshetnikov/ProveIt/tree/main/Analysis/FabiusFunction/docs/semi-formalized-research-frontiers/drafts/thue-morse/Thue_Morse_Atlas_and_Frontiers}.
The live \texttt{main} URL is not an immutable commit reference.

\bibitem{allouche-shallit}
J.-P. Allouche and J. Shallit.
The ubiquitous Prouhet--Thue--Morse sequence.
In C. Ding, T. Helleseth, and H. Niederreiter, eds.,
\emph{Sequences and Their Applications}, Springer, 1999, pp.~1--16.
\url{https://doi.org/10.1007/978-1-4471-0551-0_1}.
Author-hosted text:
\url{https://cs.uwaterloo.ca/~shallit/Papers/ubiq15.pdf}.

\bibitem{arias1982}
J. Arias de Reyna.
An infinitely differentiable function with compact support:
Definition and properties.
English translation, arXiv:1702.05442, 2017.
Original Spanish article in \emph{Revista de la Real Academia de
Ciencias Exactas, F\'isicas y Naturales de Madrid} \textbf{76}
(1982), 21--38.
\url{https://arxiv.org/abs/1702.05442}.

\bibitem{arias2018}
J. Arias de Reyna.
Arithmetic of the Fabius function.
\emph{Integers} \textbf{18} (2018), Article A51.
\url{https://math.colgate.edu/~integers/s51/s51.pdf}.
Preprint: \url{https://arxiv.org/abs/1702.06487}.

\bibitem{heyde}
C. C. Heyde.
On a property of the lognormal distribution.
\emph{Journal of the Royal Statistical Society, Series B}
\textbf{25}(2) (1963), 392--393.
\url{https://doi.org/10.1111/j.2517-6161.1963.tb00521.x}.

\bibitem{dlmf}
National Institute of Standards and Technology.
\emph{Digital Library of Mathematical Functions},
Secs.~4.22 and 24.16: infinite trigonometric products and
generalized Bernoulli polynomials.
Version 1.2.7, 15 June 2026; inspected 4 September 2026.
\url{https://dlmf.nist.gov/4.22};
\url{https://dlmf.nist.gov/24.16}.

\bibitem{cloitre}
B. Cloitre.
The Thue--Morse transform.
arXiv:2604.06243v2, 28 May 2026.
\url{https://arxiv.org/abs/2604.06243v2}.
Cited for related transform and Prouhet-partition context, not as
a source for the analytic theorems proved here.

\end{thebibliography}
\end{document}
